\documentclass[10pt, letterpaper]{article}
\usepackage[utf8]{inputenc}
\usepackage[english]{babel}
\usepackage[top=1.5cm, bottom=1.5cm, left=1.2cm, right=1.2cm]{geometry}
\usepackage{multicol}
\usepackage{listings}
\usepackage{xcolor}
\usepackage{fancyhdr}
\usepackage{amsmath}
\usepackage{amssymb}

% =================================================
% VISUAL CONFIG (Clean, Arch-style, No clutter)
% =================================================

\definecolor{keyword}{rgb}{0.0, 0.0, 0.6}
\definecolor{comment}{rgb}{0.4, 0.4, 0.4}
\definecolor{string}{rgb}{0.5, 0.0, 0.5}

\lstdefinestyle{codestyle}{
    language=C++,
    basicstyle=\ttfamily\small,       % Readable monospace font
    keywordstyle=\color{keyword}\bfseries,
    stringstyle=\color{string},
    commentstyle=\color{comment}\itshape,
    numbers=left,
    numberstyle=\tiny\color{gray},
    numbersep=5pt,
    breaklines=true,
    breakatwhitespace=true,
    tabsize=2,
    frame=b,                          % Subtle line below code blocks
    framerule=0.5pt,
    aboveskip=10pt,
    belowskip=5pt,
    showstringspaces=false
}
\lstset{style=codestyle}

% Header configuration
\pagestyle{fancy}
\fancyhf{}
\lhead{\textbf{IOI Qualifying 2026}}
\chead{Felipe Colli - "Shared Neuron" Kit}
\rhead{\thepage}
\renewcommand{\headrulewidth}{0.4pt}

% Compact Titles
\usepackage{titlesec}
\titlespacing*{\section}{0pt}{10pt}{5pt}
\titleformat{\section}{\large\bfseries}{\thesection}{1em}{}

% =================================================
% DOCUMENT CONTENT
% =================================================
\begin{document}

\begin{multicols}{2}

	% ---------------------------------------------------------
	\section{Base Template}
	% ---------------------------------------------------------
	Standard boilerplate. `long long` is life.
	\begin{lstlisting}
#include <bits/stdc++.h>
using namespace std;
typedef long long ll;
typedef vector<int> vi;
typedef pair<int,int> pii;
#define vec vector
#define pb push_back
#define all(x) x.begin(),x.end()
#define sz(x) (int)(x).size()
const int INF = 1e9;
const ll LINF = 1e18;

void fast() { ios_base::sync_with_stdio(0); cin.tie(0); }
\end{lstlisting}

	% ---------------------------------------------------------
	\section{Graph Survival Kit}
	% ---------------------------------------------------------

	\subsection{DFS (Depth First Search)}
	\textbf{Use for:} Connectivity, counting components, flood fill, checking bipartiteness.
	\textbf{Logic:} "Go deep until you can't, then backtrack."
	\begin{lstlisting}
bool vis[MAXN];
vec<int> adj[MAXN];

void dfs(int u) {
    vis[u] = true;
    // Process node 'u' here
    for (int v : adj[u]) {
        if (!vis[v]) dfs(v);
    }
}
\end{lstlisting}

	\subsection{BFS (Breadth First Search)}
	\textbf{Use for:} Shortest path in UNWEIGHTED graphs (edges = 1).
	\textbf{Logic:} "Oil spill expansion." Layer by layer.
	\begin{lstlisting}
vec<int> adj[MAXN];
int dist[MAXN]; // Initialize with -1

void bfs(int start) {
    fill(dist, dist + MAXN, -1);
    queue<int> q;

    dist[start] = 0;
    q.push(start);

    while(!q.empty()) {
        int u = q.front(); q.pop();

        for(int v : adj[u]) {
            if(dist[v] == -1) { // Not visited yet
                dist[v] = dist[u] + 1;
                q.push(v);
            }
        }
    }
}
\end{lstlisting}

	\subsection{Dijkstra (Shortest Path + Weights)}
	\textbf{Use for:} Shortest path with positive weights.
	\textbf{Logic:} It's just BFS but with a \texttt{priority\_queue}. We always process the \textit{closest} available node first.
	\begin{lstlisting}
// Adjacency list stores {neighbor, weight}
vector<pair<int, int>> adj[MAXN];
ll dist[MAXN];

void dijkstra(int start, int n) {
    // 1. Reset distances to Infinity
    for(int i=0; i<=n; i++) dist[i] = LINF;

    // 2. Min-Heap: Stores {distance, node}, smallest distance on top
    priority_queue<pair<ll,int>, vector<pair<ll,int>>, greater<pair<ll,int>>> pq;

    dist[start] = 0;
    pq.push({0, start});

    while(!pq.empty()) {
        ll d = pq.top().first;   // Current best distance
        int u = pq.top().second; // Current node
        pq.pop();

        // CRITICAL: If we found a shorter way to 'u' before extracting this
        // outdated entry from the queue, ignore it.
        if (d > dist[u]) continue;

        // Relax edges
        for (auto [v, weight] : adj[u]) {
            if (dist[u] + weight < dist[v]) {
                dist[v] = dist[u] + weight;
                pq.push({dist[v], v});
            }
        }
    }
}
\end{lstlisting}

	\subsection{DSU (Disjoint Set Union)}
	\textbf{Use for:} Grouping items, Kruskal's MST, Cycle detection.
	\textbf{Logic:} Trees of parents. `Find` gets the leader. `Unite` merges trees.
	\begin{lstlisting}
struct DSU {
    vi p, sz;
    DSU(int n) {
        p.resize(n); iota(all(p), 0); // p[i] = i
        sz.assign(n, 1);
    }
    // Path Compression: Points everything to the leader directly
    int find(int x) {
        return p[x] == x ? x : p[x] = find(p[x]);
    }
    // Union by Size: Attach small tree to big tree
    void unite(int a, int b) {
        a = find(a); b = find(b);
        if (a != b) {
            if (sz[a] < sz[b]) swap(a, b);
            p[b] = a;
            sz[a] += sz[b];
        }
    }
};
\end{lstlisting}

	% ---------------------------------------------------------
	\section{Simulation \& Search Strategies}
	% ---------------------------------------------------------

	\subsection{Two Pointers (Sliding Window)}
	\textbf{Use for:} Finding subarrays meeting a condition in $O(N)$.
	\textbf{Logic:} `R` expands the window to include elements. `L` shrinks it to maintain validity.
	\begin{lstlisting}
int l = 0, ans = 0;
for(int r = 0; r < n; r++) {
    add(a[r]); // Add new element to current state

    // While condition is broken, shrink from left
    while(!is_valid()) {
        remove(a[l]);
        l++;
    }
    // Now [l, r] is valid. Update answer.
    ans = max(ans, r - l + 1);
}
\end{lstlisting}

	\subsection{Binary Search on Answer}
	\textbf{Use for:} "Minimize the maximum X" or "Can we do it with X?".
	\textbf{Logic:} Instead of finding the best X directly, guess X and check if it's possible.
	\begin{lstlisting}
ll l = 0, r = 2e18, ans = -1;
while(l <= r) {
    ll mid = l + (r - l) / 2;
    if(check(mid)) { // check() is a simulation function
        ans = mid;
        r = mid - 1; // Try smaller (or l = mid + 1 for max)
    } else {
        l = mid + 1; // Try larger
    }
}
\end{lstlisting}

	\subsection{Coordinate Compression}
	\textbf{Use for:} Mapping huge numbers ($10^9$) to small ranks ($0 \dots N$) to fit in arrays/BITs.
	\begin{lstlisting}
vector<int> vals = v; // Copy original data
sort(all(vals));
vals.erase(unique(all(vals)), vals.end());

// Get rank (0-based index)
int get_id(int x) {
    return lower_bound(all(vals), x) - vals.begin();
}
\end{lstlisting}

	\subsection{Difference Array}
	\textbf{Use for:} Adding values to ranges $[L, R]$ efficiently.
	\begin{lstlisting}
// Add 'val' to range [l, r]
diff[l] += val;
diff[r+1] -= val;

// Reconstruct final array
for(int i=1; i<=n; i++) arr[i] = arr[i-1] + diff[i];
\end{lstlisting}

	% ---------------------------------------------------------
	\section{Math \& DP}
	% ---------------------------------------------------------

	\subsection{BinPow (Modular Exponentiation)}
	Calculates $a^b \pmod m$ in $O(\log b)$.
	\begin{lstlisting}
ll binpow(ll a, ll b, ll m) {
    a %= m; ll res = 1;
    while (b > 0) {
        if (b & 1) res = res * a % m;
        a = a * a % m;
        b >>= 1;
    }
    return res;
}
// Modular Inverse: binpow(a, MOD-2, MOD)
\end{lstlisting}

	\subsection{Combinatorics (nCr)}
	Precompute factorials first!
	\begin{lstlisting}
ll fact[MAXN], invFact[MAXN];

ll nCr(int n, int r) {
    if (r < 0 || r > n) return 0;
    return fact[n] * invFact[r] % MOD * invFact[n-r] % MOD;
}
\end{lstlisting}

	\subsection{LIS (Longest Increasing Subsequence)}
	$O(N \log N)$ approach using greedy + binary search.
	\begin{lstlisting}
vector<int> tails;
for (int x : a) {
    // lower_bound for strictly increasing
    // upper_bound for non-decreasing
    auto it = lower_bound(all(tails), x);
    if (it == tails.end()) tails.pb(x);
    else *it = x;
}
cout << tails.size();
\end{lstlisting}

\end{multicols}
\end{document}
